\documentclass[12pt]{article}
\usepackage[utf8]{inputenc}
\usepackage[T1]{fontenc}
\usepackage{hyperref}
\title{Tai Khang Language Resources for K-12 Curriculum Development}
\author{Compiled by Claude Sonnet 5 (claude-sonnet-5) for Andr\'as Kornai}
\date{}
\begin{document}
\maketitle

\section{Introduction}
Tay (Tai) Khang (ISO 639-3 tnu) is a Northern Tai (Kra-Dai) language of Laos, sometimes called just Kang. Reports on where it is spoken and by how many people differ. Wikipedia, citing Schliesinger (2003), gives about 1,500 speakers in Bolikhamsai Province. A Verbix page lists about 200 speakers. Joshua Project places the group mainly in Houaphan Province, with smaller groups elsewhere. Glottolog treats the language as endangered, with low certainty. The language is easily confused with Kh\'ang, a separate language of Vietnam, and with the code kyp, which Glottolog suggests is a duplicate entry. Very little exists for this language, and nothing educational turned up.

\section{Primary Educational Resources}
No primers, textbooks or school programmes turned up in the search results.

\section{Linguistic and Reference Materials}
\begin{itemize}
\item Wikipedia, ``Tay Khang language'': \url{https://en.wikipedia.org/wiki/Tay_Khang_language}.
\item Glottolog entry for tnu: \url{https://glottolog.org/resource/languoid/iso/tnu}.
\item Glottolog on the separate Kang entry (kyp), which is probably a duplicate: \url{https://glottolog.org/resource/languoid/iso/kyp}.
\item Glottolog entry kang1282: \url{https://glottolog.org/resource/languoid/id/kang1282}.
\item Huma-Num (CRDO) language entry: \url{https://cocoon.huma-num.fr/exist/crdo/lexvo.xql?code=tnu}. The search summary says its vitality rating is ``shifting''.
\end{itemize}

\section{Religious and Community Resources}
\begin{itemize}
\item Joshua Project profile: \url{https://joshuaproject.net/people_groups/19592}. The search summary says speakers are traditionally animist, with many having adopted Theravada Buddhism, and that no Christian believers are known. No scripture translation turned up.
\end{itemize}

\section{Digital and Audio Resources}
\begin{itemize}
\item The search summary mentions an audio recording by Michel Ferlus in the Pangloss Collection. I did not see its direct URL in the results, so I give none.
\item Verbix endangered-languages page: \url{https://docs.verbix.com/EndangeredLanguages/TayKhang}.
\end{itemize}

\section{Curriculum Development Context}
Documentation is very thin, with a very small community and likely shift toward Lao. A curriculum would first need basic documentation: a word list, an orthography decision and recorded texts. Resources for neighbouring languages should not be substituted without checking, because of the Kang/Kh\'ang confusion.

\end{document}